%% LyX 2.4.4 created this file.  For more info, see https://www.lyx.org/.
%% Do not edit unless you really know what you are doing.
\documentclass[12pt,spanish]{article}
\usepackage[utf8]{luainputenc}
\usepackage{amsmath}
\usepackage{amssymb}
\usepackage{graphicx}
\usepackage{geometry}
\geometry{verbose}
\pagestyle{headings}
\usepackage{setspace}

\makeatletter

%%%%%%%%%%%%%%%%%%%%%%%%%%%%%% LyX specific LaTeX commands.
%% Special footnote code from the package 'stblftnt.sty'
%% Author: Robin Fairbairns -- Last revised Dec 13 1996
\let\SF@@footnote\footnote
\def\footnote{\ifx\protect\@typeset@protect
    \expandafter\SF@@footnote
  \else
    \expandafter\SF@gobble@opt
  \fi
}
\expandafter\def\csname SF@gobble@opt \endcsname{\@ifnextchar[%]
  \SF@gobble@twobracket
  \@gobble
}
\edef\SF@gobble@opt{\noexpand\protect
  \expandafter\noexpand\csname SF@gobble@opt \endcsname}
\def\SF@gobble@twobracket[#1]#2{}

%%%%%%%%%%%%%%%%%%%%%%%%%%%%%% User specified LaTeX commands.
\usepackage[spanish]{babel}
\usepackage{fancyhdr}
\pagestyle{fancy}

% Limpia encabezado y pie
\fancyhf{}

% Encabezado izquierdo
\lhead{\parbox[t]{7cm}{Banco de Guatemala\\Departamento de Investigaciones Económicas\\Febrero 27, 2026}}

% Número de página al centro del pie
\cfoot{\thepage}

% Quita títulos automáticos de secciones del encabezado
\renewcommand{\sectionmark}[1]{}
\renewcommand{\subsectionmark}[1]{}

% Línea del encabezado
\renewcommand{\headrulewidth}{0.4pt}

% Sin línea en el pie
\renewcommand{\footrulewidth}{0pt}

% Espaciado
\setlength{\headsep}{24pt}
\setlength{\headheight}{36pt}

\usepackage{amsmath}
\setcounter{MaxMatrixCols}{20}

\makeatother

\usepackage{babel}
\addto\shorthandsspanish{\spanishdeactivate{~<>}}
\deactivatequoting

\begin{document}
\begin{center}
{\LARGE Banco de Guatemala}{\LARGE\par}
\par\end{center}

\begin{center}
{\LARGE Departamento de Investigaciones Económicas}{\LARGE\par}
\par\end{center}

\vspace{1cm}

\begin{center}
\includegraphics[scale=0.17]{C:/Users/Lenovo/Downloads/Logo_BG_BW}
\par\end{center}

\vspace{1cm}

\begin{onehalfspace}
\begin{center}
{\LARGE Modelo para Proyecciones Macroeconómicas}{\LARGE\par}
\par\end{center}

\begin{center}
{\large Versión 1: sin meta explicita de inflación}{\large\par}
\par\end{center}
\end{onehalfspace}

\vspace{3cm}

\begin{onehalfspace}
\begin{center}
{\large Guatemala 5 de marzo de 2026}{\large\par}
\par\end{center}
\end{onehalfspace}

\thispagestyle{empty}

\newpage

\tableofcontents{}

\thispagestyle{empty}

\lhead{Banco de Guatemala\\Departamento de Investigaciones Económicas\\Marzo 5, 2026}

\newpage
\pagenumbering{arabic}
\setcounter{page}{1}

\begin{doublespace}
\begin{center}
{\LARGE Modelo para Proyecciones Macroeconómicas}{\LARGE\par}
\par\end{center}
\end{doublespace}

A continuación se presenta el Modelo para Proyecciones Macroeconómicas
que incorpora interacciones entre las políticas monetaria y fiscal.
Para una compresión clara del mismo, la estructura del modelo se divide
en variables exógenas (externas e internas) y variables endógenas,
estas últimas clasificadas en distintos sectores:
\begin{itemize}
\item Gobierno
\item Banco Central
\item Sector Privado
\end{itemize}
Las variables que constituyen el núcleo del modelo son aquellas tildadas
con el símbolo \textasciitilde , las cuales denominamos como variables
``ola'', y cuya forma de construcción se detalla en la sección 6.
El modelo cuenta con 17 ecuaciones para las variables ``ola'' y
sus correspondientes residuos $(u)$.

\section{Variables exógenas}

\subsection{Variables exógenas externas}

\subsubsection*{Crecimiento económico externo}

\begin{equation}
\widetilde{y}_{t}^{*}=\alpha_{1,1}\widetilde{y}_{t-1}^{*}+u_{1,t}
\end{equation}

Donde $\alpha_{1,1}$ es un parámetro y $u_{1,t}$ es el residuo estocástico.\footnote{Lo mismo es cierto, de manera análoga, para las demás ecuaciones del
modelo.}

\subsubsection*{Tasa de interés externa}

\begin{equation}
\widetilde{i}_{t}^{*}=\alpha_{2,1}\widetilde{i}_{t-1}^{*}+u_{2,t}
\end{equation}


\subsubsection*{Inflación externa}

\begin{equation}
\widetilde{\pi}_{t}^{*}=\alpha_{3,1}\widetilde{\pi}_{t-1}^{*}+u_{3,t}
\end{equation}


\subsection{Variables exógenas internas}

En este modelo, la inversión pública, los ingresos totales de gobierno,
el gasto de funcionamiento del gobierno y el endeudamiento público
externo constituyen los instrumentos esenciales de la Política Fiscal.

\subsubsection*{Inversión pública}

\begin{equation}
\widetilde{\left(\frac{IPub}{PIB}\right)}_{t}=\alpha_{4,1}\widetilde{\left(\frac{IPub}{PIB}\right)}_{t-1}+u_{4,t}
\end{equation}


\subsubsection*{Ingresos totales}

\begin{equation}
\widetilde{\left(\frac{IngT}{PIB}\right)}_{t}=\alpha_{5,1}\widetilde{\left(\frac{IngT}{PIB}\right)}_{t-1}+u_{5,t}
\end{equation}


\subsubsection*{Gasto de funcionamiento}

\begin{equation}
\widetilde{\left(\frac{GFun}{PIB}\right)}_{t}=\alpha_{6,1}\widetilde{\left(\frac{GFun}{PIB}\right)}_{t-1}+u_{6,t}
\end{equation}


\subsubsection*{Deuda pública externa}

\begin{equation}
\widetilde{\left(\frac{DPE}{PIB}\right)}_{t}=\alpha_{7,1}\widetilde{\left(\frac{DPE}{PIB}\right)}_{t-1}+u_{7,t}
\end{equation}


\subsubsection*{Remesas}

Las remesas familiares que ingresan a la economía se modelan como
un proceso exógeno:

\begin{equation}
\widetilde{\left(\frac{Rem}{PIB}\right)}_{t}=\alpha_{8,1}\widetilde{\left(\frac{Rem}{PIB}\right)}_{t-1}+u_{8,t}
\end{equation}


\section{Variables endógenas}

Las variables endógenas se clasifican según su sector dentro de la
economía doméstica.

\subsection{Gobierno}

\subsubsection*{Deuda pública interna\protect\footnote{La ecuación (9), que caracteriza al cociente de la deuda pública interna
a PIB, es una ecuación de comportamiento. Su compatibilidad con la
definición contable de la deuda pública se aborda en la sección 4.}}

\begin{equation}
\widetilde{\left(\frac{DPI}{PIB}\right)}_{t}=\alpha_{9,1}\widetilde{\left(\frac{DPI}{PIB}\right)}_{t-1}+\alpha_{9,2}\widetilde{\left(\frac{DPE}{PIB}\right)}_{t-1}+\alpha_{9,3}\widetilde{\left(\frac{Def}{PIB}\right)}_{t-1}+u_{9,t}
\end{equation}


\subsection{Banco Central}

\subsubsection*{Activos externos netos}

\begin{equation}
\widetilde{\left(\frac{AEN}{PIB}\right)}_{t}=\alpha_{10,1}\widetilde{\left(\frac{AEN}{PIB}\right)}_{t-1}+\alpha_{10,2}\widetilde{\left(\frac{DPE}{PIB}\right)}_{t-1}+\alpha_{10,3}\widetilde{\pi}_{t-1}^{*}+\alpha_{10,4}\widetilde{e}_{t-1}+\alpha_{10,5}\widetilde{\pi}_{t-1}+u_{10,t}
\end{equation}

Dependen de su rezago y de los rezagos de las variables siguientes:
\begin{itemize}
\item $\widetilde{\left(\frac{DPE}{PIB}\right)}$= Deuda pública externa
\item $\widetilde{\pi}^{*}$ = Inflación externa
\item $\widetilde{e}$ = Depreciación nominal
\item $\widetilde{\pi}$ = Inflación doméstica
\end{itemize}
\newpage

\subsubsection*{Activos internos netos}

\begin{equation}
\widetilde{\left(\frac{AIN}{PIB}\right)}_{t}=\alpha_{11,1}\widetilde{\left(\frac{AIN}{PIB}\right)}_{t-1}+\alpha_{11,2}\widetilde{\left(\frac{DPI}{PIB}\right)}_{t-1}+\alpha_{11,3}\widetilde{i}_{t-1}+\alpha_{11,4}\widetilde{\pi}_{t-1}+u_{11,t}
\end{equation}

Dependen su rezago y de los rezagos de las variables siguientes:
\begin{itemize}
\item $\widetilde{\left(\frac{DPI}{PIB}\right)}$= Deuda pública interna
\item $\widetilde{i}$ = Tasa de interés líder de política monetaria
\item $\widetilde{\pi}$ = Inflación doméstica
\end{itemize}

\subsubsection*{Tasa de interés líder de política monetaria (regla de política monetaria)}

\begin{equation}
\widetilde{i}_{t}=\alpha_{12,1}\widetilde{i}_{t-1}+\alpha_{12,2}\widetilde{i}_{t-1}^{*}+\alpha_{12,3}\widetilde{\pi}_{t-1}+\alpha_{12,4}\widetilde{y}_{t-1}+u_{12,t}
\end{equation}

Depende de su rezago y de los rezagos de las variables siguientes:
\begin{itemize}
\item $\widetilde{i}^{*}$ = Tasa de interés externa
\item $\widetilde{\pi}$ = Inflación doméstica
\item $\widetilde{y}$ = Crecimiento económico doméstico 
\end{itemize}

\subsection{Sector privado}

\subsubsection*{Depreciación nominal}

\begin{equation}
\widetilde{e}_{t}=\alpha_{13,1}\widetilde{e}_{t-1}+\alpha_{13,2}\widetilde{\left(\frac{DPEC}{PIB}\right)}_{t-1}+\alpha_{13,3}\left(\widetilde{i}_{t-1}-\widetilde{i}_{t-1}^{*}-\widetilde{\rho}_{t-1}\right)+u_{13,t}
\end{equation}

Depende de su rezago y de los rezagos de las variables siguientes:
\begin{itemize}
\item $\widetilde{\left(\frac{DPEC}{PIB}\right)}$= Deuda pública externa
consolidada
\item $\widetilde{i}$ = Tasa de interés líder de política monetaria
\item $\widetilde{i}^{*}$ = Tasa de interés externa
\item $\widetilde{\rho}$ = Prima por riesgo país
\end{itemize}

\subsubsection*{Prima por Riesgo-País}

\begin{equation}
\widetilde{\rho}_{t}=\alpha_{14,1}\widetilde{\rho}_{t-1}+\alpha_{14,2}\widetilde{\left(\frac{DPTC}{PIB}\right)}_{t-1}+u_{14,t}
\end{equation}

Depende de su rezago y del rezagos de la siguiente variable:
\begin{itemize}
\item $\widetilde{\left(\frac{DPTC}{PIB}\right)}$= Deuda pública total
consolidada
\end{itemize}

\subsubsection*{Inflación (curva de Phillips)}

\begin{equation}
\widetilde{\pi}_{t}=\alpha_{15,1}\widetilde{\pi}_{t-1}+\alpha_{15,2}\widetilde{y}_{t-1}+\alpha_{15,3}\widetilde{\pi}_{t-1}^{*}+\alpha_{15,4}\widetilde{e}_{t-1}+\alpha_{15,5}\widetilde{\left(\frac{MO}{PIB}\right)}_{t-1}+u_{15,t}
\end{equation}

Depende de su rezago y:
\begin{itemize}
\item $\widetilde{y}$ = Crecimiento económico doméstico
\item $\widetilde{\pi}^{*}$ = Tasa de inflación externa
\item $\widetilde{e}$ = Depreciación nominal
\item $\widetilde{\left(\frac{MO}{PIB}\right)}$ = Base monetaria
\end{itemize}

\subsubsection*{Crecimiento económico (curva IS)}

\begin{align*}
\widetilde{y}_{t} & =\alpha_{16,1}\widetilde{y}_{t-1}+\alpha_{16,2}\widetilde{y}_{t-1}^{*}+\alpha_{16,3}\widetilde{\left(\frac{Rem}{PIB}\right)}_{t-1}
\end{align*}

\begin{equation}
+\alpha_{16,4}\widetilde{\left(\frac{ITotal}{PIB}\right)}_{t-1}+\alpha_{16,5}\widetilde{\left(\frac{MO}{PIB}\right)}_{t-1}+u_{16,t}
\end{equation}

Depende de su rezago y de los rezagos de las variables siguientes:
\begin{itemize}
\item $\widetilde{y}^{*}$ = Crecimiento económico externo
\item $\widetilde{\left(\frac{ITotal}{PIB}\right)}$ = Inversión total
\item $\widetilde{\left(\frac{MO}{PIB}\right)}$ = Base monetaria
\end{itemize}

\subsubsection*{Inversión Privada}

\begin{equation}
\widetilde{\left(\frac{IPriv}{PIB}\right)}_{t}=\alpha_{17,1}\widetilde{\left(\frac{IPriv}{PIB}\right)}_{t-1}+\alpha_{17,2}\widetilde{i}_{t-1}+\alpha_{17,3}\widetilde{\left(\frac{DPIC}{PIB}\right)}_{t-1}+\alpha_{17,4}\widetilde{\rho}_{t-1}+u_{17,t}
\end{equation}

Depende de su rezago y de los rezagos de las variables siguientes:
\begin{itemize}
\item $\widetilde{i}$ = Tasa de interés Líder de Política Monetaria
\item $\widetilde{\left(\frac{DPIC}{PIB}\right)}$ = Deuda pública interna
consolidada
\item $\tilde{\rho}$ = Prima por riesgo país
\end{itemize}

\section{Definiciones}

A continuación se presentan un conjunto de definiciones separadas
del núcleo del modelo.

\subsection*{Déficit fiscal}

\begin{equation}
\widetilde{\left(\frac{Def}{PIB}\right)}_{t}=\widetilde{\left(\frac{GFun}{PIB}\right)}_{t}+\widetilde{\left(\frac{IPub}{PIB}\right)}_{t}-\widetilde{\left(\frac{IngT}{PIB}\right)}_{t}
\end{equation}


\subsection*{Deuda pública total\protect\footnote{La ecuación (19) representa la ley de movimiento de la deuda pública
total, en tanto que la ecuación (20) descompone la deuda pública total
en un componente externo y otro interno. La compatibilidad de ambas
definiciones se explica en la sección 4.}}

\begin{equation}
\widetilde{\left(\frac{DPT}{PIB}\right)}_{t}=\widetilde{\left(\frac{DPT}{PIB}\right)}_{t-1}+\widetilde{\left(\frac{Def}{PIB}\right)}_{t}
\end{equation}

\begin{equation}
\widetilde{\left(\frac{DPT}{PIB}\right)}_{t}=\widetilde{\left(\frac{DPE}{PIB}\right)}_{t}+\widetilde{\left(\frac{DPI}{PIB}\right)}_{t}
\end{equation}


\subsection*{Base monetaria}

\begin{equation}
\widetilde{\left(\frac{MO}{PIB}\right)}_{t}=\widetilde{\left(\frac{AEN}{PIB}\right)}_{t}+\widetilde{\left(\frac{AIN}{PIB}\right)}_{t}
\end{equation}


\subsection*{Inversión Total}

\begin{equation}
\widetilde{\left(\frac{ITotal}{PIB}\right)}_{t}=\widetilde{\left(\frac{IPub}{PIB}\right)}_{t}+\widetilde{\left(\frac{IPriv}{PIB}\right)}_{t}
\end{equation}


\subsection*{Deuda pública interna consolidada}

\begin{equation}
\widetilde{\left(\frac{DPIC}{PIB}\right)}_{t}=\widetilde{\left(\frac{DPI}{PIB}\right)}_{t}-\widetilde{\left(\frac{AIN}{PIB}\right)}_{t}
\end{equation}


\subsection*{Deuda pública externa consolidada}

\begin{equation}
\widetilde{\left(\frac{DPEC}{PIB}\right)}_{t}=\widetilde{\left(\frac{DPE}{PIB}\right)}_{t}-\widetilde{\left(\frac{AEN}{PIB}\right)}_{t}
\end{equation}


\subsection*{Deuda pública total consolidada}

\begin{equation}
\widetilde{\left(\frac{DPTC}{PIB}\right)}_{t}=\widetilde{\left(\frac{DPEC}{PIB}\right)}_{t}+\widetilde{\left(\frac{DPIC}{PIB}\right)}_{t}
\end{equation}


\section{Sustitución de variables}

\subsection*{Déficit fiscal en (9)}

Sustituyendo la ecuación (18), rezagada un período, en la ecuación
(9), se obtiene:

\[
\widetilde{\left(\frac{DPI}{PIB}\right)}_{t}=\alpha_{9,1}\widetilde{\left(\frac{DPI}{PIB}\right)}_{t-1}+\alpha_{9,2}\widetilde{\left(\frac{DPE}{PIB}\right)}_{t-1}
\]

\begin{equation}
+\alpha_{9,3}\left[\widetilde{\left(\frac{GFun}{PIB}\right)}_{t-1}+\widetilde{\left(\frac{IPub}{PIB}\right)}_{t-1}-\widetilde{\left(\frac{IngT}{PIB}\right)}_{t-1}\right]+u_{9,t}
\end{equation}


\subsection*{Deuda púbica externa consolidada en (13)}

Sustituyendo la ecuación (24), rezagada un período, en la ecuación
(13), resulta:

\begin{equation}
\widetilde{e}_{t}=\alpha_{13,1}\widetilde{e}_{t-1}+\alpha_{13,2}\left[\widetilde{\left(\frac{DPE}{PIB}\right)}_{t-1}-\widetilde{\left(\frac{AEN}{PIB}\right)}_{t-1}\right]+\alpha_{13,3}\left(\widetilde{i}_{t-1}-\widetilde{i}_{t-1}^{*}-\widetilde{\rho}_{t-1}\right)+u_{13,t}
\end{equation}


\subsection*{Deuda pública total consolidada en (14)}

Sustituyendo\footnote{Secuencialmente y con los rezagos del caso.}
las ecuaciones (25) (24) y (23) en la ecuación (14), se obtiene:

\begin{equation}
\widetilde{\rho}_{t}=\alpha_{14,1}\widetilde{\rho}_{t-1}+\alpha_{14,2}\left[\widetilde{\left(\frac{DPEC}{PIB}\right)}_{t-1}+\widetilde{\left(\frac{DPIC}{PIB}\right)}_{t-1}\right]+u_{14,t}
\end{equation}

luego:

\begin{equation}
\widetilde{\rho}_{t}=\alpha_{14,1}\widetilde{\rho}_{t-1}+\alpha_{14,2}\left[\widetilde{\left(\frac{DPE}{PIB}\right)}_{t-1}-\widetilde{\left(\frac{AEN}{PIB}\right)}_{t-1}+\widetilde{\left(\frac{DPI}{PIB}\right)}_{t-1}-\widetilde{\left(\frac{AIN}{PIB}\right)}_{t-1}\right]+u_{14,t}
\end{equation}


\subsection*{Base Monetaria en (15)}

Sustituyendo la ecuación (21) en la ecuación (15), se obtiene:

\begin{equation}
\widetilde{\pi}_{t}=\alpha_{15,1}\widetilde{\pi}_{t-1}+\alpha_{15,2}\widetilde{y}_{t-1}+\alpha_{15,3}\widetilde{\pi}_{t-1}^{*}+\alpha_{15,4}\widetilde{e}_{t-1}+\alpha_{15,5}\left[\widetilde{\left(\frac{AEN}{PIB}\right)}_{t-1}+\widetilde{\left(\frac{AIN}{PIB}\right)}_{t-1}\right]+u_{15,t}
\end{equation}

\newpage

\subsection*{Inversión total y base monetaria en (16)}

Sustituyendo las ecuaciones (22) y (21) en la ecuación (16), obtenemos:

\begin{align*}
\widetilde{y}_{t} & =\alpha_{16,1}\widetilde{y}_{t-1}+\alpha_{16,2}\widetilde{y}_{t-1}^{*}+\alpha_{16,3}\widetilde{\left(\frac{Rem}{PIB}\right)}_{t-1}
\end{align*}

\begin{equation}
+\alpha_{16,4}\left[\widetilde{\left(\frac{IPubl}{PIB}\right)}_{t-1}+\widetilde{\left(\frac{IPriv}{PIB}\right)}_{t-1}\right]+\alpha_{16,5}\left[\widetilde{\left(\frac{AEN}{PIB}\right)}_{t-1}+\widetilde{\left(\frac{AIN}{PIB}\right)}_{t-1}\right]+u_{16,t}
\end{equation}


\subsection*{Deuda pública interna consolidada en (17)}

Sustituyendo la ecuación (23), rezagada un período, en la ecuación
(17), resulta:

\[
\widetilde{\left(\frac{IPriv}{PIB}\right)}_{t}=\alpha_{17,1}\widetilde{\left(\frac{IPriv}{PIB}\right)}_{t-1}+\alpha_{17,2}\widetilde{i}_{t-1}
\]

\begin{equation}
+\alpha_{17,3}\left[\widetilde{\left(\frac{DPI}{PIB}\right)}_{t-1}-\widetilde{\left(\frac{AIN}{PIB}\right)}_{t-1}\right]+\alpha_{17,4}\widetilde{\rho}_{t-1}+u_{17,t}
\end{equation}


\subsection*{Deuda pública total en (19)}

Sustituyendo la ecuación (20) en la ecuación (19), con y sin rezago,
obtenemos:

\[
\widetilde{\left(\frac{DPE}{PIB}\right)}_{t}+\widetilde{\left(\frac{DPI}{PIB}\right)}_{t}=\widetilde{\left(\frac{DPE}{PIB}\right)}_{t-1}+\widetilde{\left(\frac{DPI}{PIB}\right)}_{t-1}+\widetilde{\left(\frac{Def}{PIB}\right)}_{t}
\]

\begin{equation}
\Longrightarrow\widetilde{\left(\frac{DPI}{PIB}\right)}_{t}=\widetilde{\left(\frac{DPE}{PIB}\right)}_{t-1}+\widetilde{\left(\frac{DPI}{PIB}\right)}_{t-1}+\widetilde{\left(\frac{Def}{PIB}\right)}_{t}-\widetilde{\left(\frac{DPE}{PIB}\right)}_{t}
\end{equation}

\newpage

Luego, sustituyendo la ecuación (9) en la ecuación (33), se tiene:

\[
\alpha_{9,1}\widetilde{\left(\frac{DPI}{PIB}\right)}_{t-1}+\alpha_{9,2}\widetilde{\left(\frac{DPE}{PIB}\right)}_{t-1}+\alpha_{9,3}\widetilde{\left(\frac{Def}{PIB}\right)}_{t-1}+u_{9,t}=\widetilde{\left(\frac{DPE}{PIB}\right)}_{t-1}+\widetilde{\left(\frac{DPI}{PIB}\right)}_{t-1}
\]

\[
-\widetilde{\left(\frac{DPE}{PIB}\right)}_{t}+\widetilde{\left(\frac{Def}{PIB}\right)}_{t}
\]

Despejando para $u_{9,t}$ y reagrupando:

\[
\Longrightarrow u_{9,t}=\left(1-\alpha_{9,1}\right)\widetilde{\left(\frac{DPI}{PIB}\right)}_{t-1}+\left(1-\alpha_{9,2}\right)\widetilde{\left(\frac{DPE}{PIB}\right)}_{t-1}
\]

\begin{equation}
-\alpha_{9,3}\widetilde{\left(\frac{Def}{PIB}\right)}_{t-1}-\widetilde{\left(\frac{DPE}{PIB}\right)}_{t}+\widetilde{\left(\frac{Def}{PIB}\right)}_{t}
\end{equation}

Ahora, sustituyendo (7) en (34), se obtiene:

\[
u_{9,t}=\left(1-\alpha_{9,1}\right)\widetilde{\left(\frac{DPI}{PIB}\right)}_{t-1}+\left(1-\alpha_{9,2}\right)\widetilde{\left(\frac{DPE}{PIB}\right)}_{t-1}
\]

\[
-\alpha_{9,3}\widetilde{\left(\frac{Def}{PIB}\right)}_{t-1}-\left[\alpha_{7,1}\widetilde{\left(\frac{DPE}{PIB}\right)}_{t-1}+u_{7,t}\right]+\widetilde{\left(\frac{Def}{PIB}\right)}_{t}
\]

\[
u_{9,t}=\left(1-\alpha_{9,1}\right)\widetilde{\left(\frac{DPI}{PIB}\right)}_{t-1}+\left(1-\alpha_{9,2}\right)\widetilde{\left(\frac{DPE}{PIB}\right)}_{t-1}
\]

\[
-\alpha_{9,3}\widetilde{\left(\frac{Def}{PIB}\right)}_{t-1}-\alpha_{7,1}\widetilde{\left(\frac{DPE}{PIB}\right)}_{t-1}-u_{7,t}+\widetilde{\left(\frac{Def}{PIB}\right)}_{t}
\]

\begin{equation}
\Longrightarrow u_{9,t}=\left(1-\alpha_{9,1}\right)\widetilde{\left(\frac{DPI}{PIB}\right)}_{t-1}+\left(1-\alpha_{7,1}-\alpha_{9,2}\right)\widetilde{\left(\frac{DPE}{PIB}\right)}_{t-1}-\alpha_{9,3}\widetilde{\left(\frac{Def}{PIB}\right)}_{t-1}+\widetilde{\left(\frac{Def}{PIB}\right)}_{t}-u_{7,t}
\end{equation}

\newpage

Seguidamente, se sustituye la ecuación (18), tanto de manera rezagada
como en forma contemporánea, en la ecuación (35):

\[
u_{9,t}=\left(1-\alpha_{9,1}\right)\widetilde{\left(\frac{DPI}{PIB}\right)}_{t-1}+\left(1-\alpha_{7,1}-\alpha_{9,2}\right)\widetilde{\left(\frac{DPE}{PIB}\right)}_{t-1}
\]

\[
-\alpha_{9,3}\left[\widetilde{\left(\frac{GFun}{PIB}\right)}_{t-1}+\widetilde{\left(\frac{IPub}{PIB}\right)}_{t-1}-\widetilde{\left(\frac{IngT}{PIB}\right)}_{t-1}\right]
\]

\begin{equation}
+\left[\widetilde{\left(\frac{GFun}{PIB}\right)}_{t}+\widetilde{\left(\frac{IPub}{PIB}\right)}_{t}-\widetilde{\left(\frac{IngT}{PIB}\right)}_{t}\right]-u_{7,t}
\end{equation}

Ahora procedemos a sustituir las ecuaciones (4), (5) y (6) en la ecuación
(36):

\[
u_{9,t}=\left(1-\alpha_{9,1}\right)\widetilde{\left(\frac{DPI}{PIB}\right)}_{t-1}+\left(1-\alpha_{7,1}-\alpha_{9,2}\right)\widetilde{\left(\frac{DPE}{PIB}\right)}_{t-1}
\]

\[
-\alpha_{9,3}\left[\widetilde{\left(\frac{GFun}{PIB}\right)}_{t-1}+\widetilde{\left(\frac{IPub}{PIB}\right)}_{t-1}-\widetilde{\left(\frac{IngT}{PIB}\right)}_{t-1}\right]
\]

\[
+\left\{ \left[\alpha_{6,1}\widetilde{\left(\frac{GFun}{PIB}\right)}_{t-1}+u_{6,t}\right]+\left[\alpha_{4,1}\widetilde{\left(\frac{IPub}{PIB}\right)}_{t-1}+u_{4,t}\right]-\left[\alpha_{5,1}\widetilde{\left(\frac{IngT}{PIB}\right)}_{t-1}+u_{5,t}\right]\right\} -u_{7,1}
\]

\[
\Longrightarrow u_{9,t}=\left(1-\alpha_{9,1}\right)\widetilde{\left(\frac{DPI}{PIB}\right)}_{t-1}+\left(1-\alpha_{7,1}-\alpha_{9,2}\right)\widetilde{\left(\frac{DPE}{PIB}\right)}_{t-1}+\left(\alpha_{6,1}-\alpha_{9,3}\right)\widetilde{\left(\frac{GFun}{PIB}\right)}_{t-1}
\]

\begin{equation}
+\left(\alpha_{4,1}-\alpha_{9,3}\right)\widetilde{\left(\frac{IPub}{PIB}\right)}_{t-1}-\left(\alpha_{5,1}+\alpha_{9,3}\right)\widetilde{\left(\frac{IngT}{PIB}\right)}_{t-1}+u_{4,t}+u_{6,t}-u_{5,t}-u_{7,t}
\end{equation}


\section{Presentación compacta del modelo}

El modelo para Proyecciones Macroeconómicas puede reescribirse, de
manera más compacta (aunque menos intuitiva), mediante las 18 ecuaciones
siguientes:

\noindent (1), (2), (3), (4), (5), (6), (7), (8), (26), (10), (11),
(12), (27), (29), (30), (31), (32) y (37)

\noindent A continuación se reproducen las recién referidas ecuaciones,
pero esta vez de manera consecutiva y con nueva numeración:

\begin{equation}
\widetilde{y}_{t}^{*}=\alpha_{1,1}\widetilde{y}_{t-1}^{*}+u_{1,t}
\end{equation}

\begin{equation}
\widetilde{i}_{t}^{*}=\alpha_{2,1}\widetilde{i}_{t-1}^{*}+u_{2,t}
\end{equation}

\begin{equation}
\widetilde{\pi}_{t}^{*}=\alpha_{3,1}\widetilde{\pi}_{t-1}^{*}+u_{3,t}
\end{equation}

\begin{equation}
\widetilde{\left(\frac{IPub}{PIB}\right)}_{t}=\alpha_{4,1}\widetilde{\left(\frac{IPub}{PIB}\right)}_{t-1}+u_{4,t}
\end{equation}

\begin{equation}
\widetilde{\left(\frac{IngT}{PIB}\right)}_{t}=\alpha_{5,1}\widetilde{\left(\frac{IngT}{PIB}\right)}_{t-1}+u_{5,t}
\end{equation}

\begin{equation}
\widetilde{\left(\frac{GFun}{PIB}\right)}_{t}=\alpha_{6,1}\widetilde{\left(\frac{GFun}{PIB}\right)}_{t-1}+u_{6,t}
\end{equation}

\begin{equation}
\widetilde{\left(\frac{DPE}{PIB}\right)}_{t}=\alpha_{7,1}\widetilde{\left(\frac{DPE}{PIB}\right)}_{t-1}+u_{7,t}
\end{equation}

\begin{equation}
\widetilde{\left(\frac{Rem}{PIB}\right)}_{t}=\alpha_{8,1}\widetilde{\left(\frac{Rem}{PIB}\right)}_{t-1}+u_{8,t}
\end{equation}

\[
\widetilde{\left(\frac{DPI}{PIB}\right)}_{t}=\alpha_{9,1}\widetilde{\left(\frac{DPI}{PIB}\right)}_{t-1}+\alpha_{9,2}\widetilde{\left(\frac{DPE}{PIB}\right)}_{t-1}
\]

\begin{equation}
+\alpha_{9,3}\left[\widetilde{\left(\frac{GFun}{PIB}\right)}_{t-1}+\widetilde{\left(\frac{IPub}{PIB}\right)}_{t-1}-\widetilde{\left(\frac{IngT}{PIB}\right)}_{t-1}\right]+u_{9,t}
\end{equation}

\begin{equation}
\widetilde{\left(\frac{AEN}{PIB}\right)}_{t}=\alpha_{10,1}\widetilde{\left(\frac{AEN}{PIB}\right)}_{t-1}+\alpha_{10,2}\widetilde{\left(\frac{DPE}{PIB}\right)}_{t-1}+\alpha_{10,3}\widetilde{\pi}_{t-1}^{*}+\alpha_{10,4}\widetilde{e}_{t-1}+\alpha_{10,5}\widetilde{\pi}_{t-1}+u_{10,t}
\end{equation}

\begin{equation}
\widetilde{\left(\frac{AIN}{PIB}\right)}_{t}=\alpha_{11,1}\widetilde{\left(\frac{AIN}{PIB}\right)}_{t-1}+\alpha_{11,2}\widetilde{\left(\frac{DPI}{PIB}\right)}_{t-1}+\alpha_{11,3}\widetilde{i}_{t-1}+\alpha_{11,4}\widetilde{\pi}_{t-1}+u_{11,t}
\end{equation}

\begin{equation}
\widetilde{i}_{t}=\alpha_{12,1}\widetilde{i}_{t-1}+\alpha_{12,2}\widetilde{i}_{t-1}^{*}+\alpha_{12,3}\widetilde{\pi}_{t-1}+\alpha_{12,4}\widetilde{y}_{t-1}+u_{12,t}
\end{equation}

\begin{equation}
\widetilde{e}_{t}=\alpha_{13,1}\widetilde{e}_{t-1}+\alpha_{13,2}\left[\widetilde{\left(\frac{DPE}{PIB}\right)}_{t-1}-\widetilde{\left(\frac{AEN}{PIB}\right)}_{t-1}\right]+\alpha_{13,3}\left(\widetilde{i}_{t-1}-\widetilde{i}_{t-1}^{*}-\widetilde{\rho}_{t-1}\right)+u_{13,t}
\end{equation}

\begin{equation}
\widetilde{\rho}_{t}=\alpha_{14,1}\widetilde{\rho}_{t-1}+\alpha_{14,2}\left[\widetilde{\left(\frac{DPI}{PIB}\right)}_{t-1}-\widetilde{\left(\frac{AIN}{PIB}\right)}_{t-1}+\widetilde{\left(\frac{DPE}{PIB}\right)}_{t-1}-\widetilde{\left(\frac{AEN}{PIB}\right)}_{t-1}\right]+u_{14,t}
\end{equation}

\begin{equation}
\widetilde{\pi}_{t}=\alpha_{15,1}\widetilde{\pi}_{t-1}+\alpha_{15,2}\widetilde{y}_{t-1}+\alpha_{15,3}\widetilde{\pi}_{t-1}^{*}+\alpha_{15,4}\widetilde{e}_{t-1}+\alpha_{15,5}\left[\widetilde{\left(\frac{AEN}{PIB}\right)}_{t-1}+\widetilde{\left(\frac{AIN}{PIB}\right)}_{t-1}\right]+u_{15,t}
\end{equation}

\begin{align*}
\widetilde{y}_{t} & =\alpha_{16,1}\widetilde{y}_{t-1}+\alpha_{16,2}\widetilde{y}_{t-1}^{*}+\alpha_{16,3}\widetilde{\left(\frac{Rem}{PIB}\right)}_{t-1}
\end{align*}

\begin{equation}
+\alpha_{16,4}\left[\widetilde{\left(\frac{IPubl}{PIB}\right)}_{t-1}+\widetilde{\left(\frac{IPriv}{PIB}\right)}_{t-1}\right]+\alpha_{16,5}\left[\widetilde{\left(\frac{AEN}{PIB}\right)}_{t-1}+\widetilde{\left(\frac{AIN}{PIB}\right)}_{t-1}\right]+u_{16,t}
\end{equation}

\begin{equation}
\widetilde{\left(\frac{IPriv}{PIB}\right)}_{t}=\alpha_{17,1}\widetilde{\left(\frac{IPriv}{PIB}\right)}_{t-1}+\alpha_{17,2}\widetilde{i}_{t-1}+\alpha_{17,3}\left[\widetilde{\left(\frac{DPI}{PIB}\right)}_{t-1}-\widetilde{\left(\frac{AIN}{PIB}\right)}_{t-1}\right]+\alpha_{17,4}\widetilde{\rho}_{t-1}+u_{17,t}
\end{equation}

\[
u_{9,t}=\left(1-\alpha_{9,1}\right)\widetilde{\left(\frac{DPI}{PIB}\right)}_{t-1}+\left(1-\alpha_{7,1}-\alpha_{9,2}\right)\widetilde{\left(\frac{DPE}{PIB}\right)}_{t-1}+\left(\alpha_{6,1}-\alpha_{9,3}\right)\widetilde{\left(\frac{GFun}{PIB}\right)}_{t-1}
\]

\begin{equation}
+\left(\alpha_{4,1}-\alpha_{9,3}\right)\widetilde{\left(\frac{IPub}{PIB}\right)}_{t-1}-\left(\alpha_{5,1}+\alpha_{9,3}\right)\widetilde{\left(\frac{IngT}{PIB}\right)}_{t-1}+u_{4,t}+u_{6,t}-u_{5,t}-u_{7,t}
\end{equation}


\section{Construcción de las variables ``ola'' y recuperación de Variables
Observables}

Una vez expuestas las ecuaciones fundamentales del modelo, el siguiente
paso es definir cómo están construidas las variables que lo componen.
La estructura general de su construcción se detalla a continuación.

\subsection*{Definición de variables ``ola''}

Las variables ``ola'' se construyen como la diferencia entre la
variable observable en el periodo $t$ menos el rezago de un promedio
móvil $T$ periodos. Estas variables son las que forman parte del
modelo propiamente tal que:

\begin{equation}
\widetilde{x}_{t}=x_{t}-\overline{x}_{t-1}
\end{equation}


\subsection*{Recuperación de variables observables}

Por su parte, las variables observables son una reconstrucción a partir
de los pronósticos de las variables ``ola'' y constituyen las variables
originales de la economía representada por el modelo. Se definen como
la suma de la correspondiente variable ``ola'' en el periodo $t$
y el respectivo rezago de un promedio móvil de $T$ periodos:

\begin{equation}
x_{t}=\widetilde{x}_{t}+\overline{x}_{t-1}
\end{equation}


\subsection*{Promedio móvil}

El promedio móvil que aparece en las ecuaciones $(56)$ y $(57)$
se define de la siguiente manera:

\begin{equation}
\overline{x}_{t-1}=\frac{1}{T}\sum_{j=t-T}^{t-1}x_{j}
\end{equation}

\noindent A continuación, se aplican estas definiciones a todas las
variables del modelo:

\newpage

\subsection{Variables exógenas externas}

\subsubsection*{Crecimiento económico externo}

\begin{equation}
y_{t}^{*}=\widetilde{y}_{t}^{*}+\overline{y}_{t-1}^{*}
\end{equation}

\begin{equation}
\widetilde{y}_{t}^{*}=y_{t}^{*}-\overline{y}_{t-1}^{*}
\end{equation}

\begin{equation}
\overline{y}_{t-1}^{*}=\frac{1}{T}\sum_{j=t-T}^{t-1}y_{j}^{*}
\end{equation}


\subsubsection*{Tasa de interés externa}

\begin{equation}
i_{t}^{*}=\widetilde{i}_{t}^{*}+\overline{i}_{t-1}^{*}
\end{equation}

\begin{equation}
\widetilde{i}_{t}^{*}=i_{t}^{*}-\overline{i}_{t-1}^{*}
\end{equation}

\begin{equation}
\overline{i}_{t-1}^{*}=\frac{1}{T}\sum_{j=t-T}^{t-1}i_{j}^{*}
\end{equation}


\subsubsection*{Inflación externa}

\begin{equation}
\pi_{t}^{*}=\widetilde{\pi}_{t}^{*}+\overline{\pi}_{t-1}^{*}
\end{equation}

\begin{equation}
\widetilde{\pi}_{t}^{*}=\pi_{t}^{*}-\overline{\pi}_{t-1}^{*}
\end{equation}

\begin{equation}
\overline{\pi}_{t-1}^{*}=\frac{1}{T}\sum_{j=t-T}^{t-1}\pi_{j}^{*}
\end{equation}


\subsection{Variables exógenas internas}

\subsubsection*{Inversión Pública}

\begin{equation}
\left(\frac{IPub}{PIB}\right)_{t}=\widetilde{\left(\frac{IPub}{PIB}\right)}_{t}+\overline{\left(\frac{IPub}{PIB}\right)}_{t-1}
\end{equation}

\begin{equation}
\widetilde{\left(\frac{IPub}{PIB}\right)}_{t}=\left(\frac{IPub}{PIB}\right)_{t}-\overline{\left(\frac{IPub}{PIB}\right)}_{t-1}
\end{equation}

\begin{equation}
\overline{\left(\frac{IPub}{PIB}\right)}_{t-1}=\frac{1}{T}\sum_{j=t-T}^{t-1}\left(\frac{IPub}{PIB}\right)_{j}
\end{equation}


\subsubsection*{Ingresos totales}

\begin{equation}
\left(\frac{IngT}{PIB}\right)_{t}=\widetilde{\left(\frac{IngT}{PIB}\right)}_{t}+\overline{\left(\frac{IngT}{PIB}\right)}_{t-1}
\end{equation}

\begin{equation}
\widetilde{\left(\frac{IngT}{PIB}\right)}_{t}=\left(\frac{IngT}{PIB}\right)_{t}-\overline{\left(\frac{IngT}{PIB}\right)}_{t-1}
\end{equation}

\begin{equation}
\overline{\left(\frac{IngT}{PIB}\right)}_{t-1}=\frac{1}{T}\sum_{j=t-T}^{t-1}\left(\frac{IngT}{PIB}\right)_{j}
\end{equation}


\subsubsection*{Gasto en funcionamiento del Gobierno}

\begin{equation}
\left(\frac{GFun}{PIB}\right)_{t}=\widetilde{\left(\frac{GFun}{PIB}\right)}_{t}+\overline{\left(\frac{GFun}{PIB}\right)}_{t-1}
\end{equation}

\begin{equation}
\widetilde{\left(\frac{GFun}{PIB}\right)}_{t}=\left(\frac{GFun}{PIB}\right)_{t}-\overline{\left(\frac{GFun}{PIB}\right)}_{t-1}
\end{equation}

\begin{equation}
\overline{\left(\frac{GFun}{PIB}\right)}_{t-1}=\frac{1}{T}\sum_{j=t-T}^{t-1}\left(\frac{GFun}{PIB}\right)_{j}
\end{equation}


\subsubsection*{Deuda pública externa}

\begin{equation}
\left(\frac{DPE}{PIB}\right)_{t}=\widetilde{\left(\frac{DPE}{PIB}\right)}_{t}+\overline{\left(\frac{DPE}{PIB}\right)}_{t-1}
\end{equation}

\begin{equation}
\widetilde{\left(\frac{DPE}{PIB}\right)}_{t}=\left(\frac{DPE}{PIB}\right)_{t}-\overline{\left(\frac{DPE}{PIB}\right)}_{t-1}
\end{equation}

\begin{equation}
\overline{\left(\frac{DPE}{PIB}\right)}_{t-1}=\frac{1}{T}\sum_{j=t-T}^{t-1}\left(\frac{DPE}{PIB}\right)_{j}
\end{equation}


\subsubsection*{Remesas}

\begin{equation}
\left(\frac{Rem}{PIB}\right)_{t}=\widetilde{\left(\frac{Rem}{PIB}\right)}_{t}+\overline{\left(\frac{Rem}{PIB}\right)}_{t-1}
\end{equation}

\begin{equation}
\widetilde{\left(\frac{Rem}{PIB}\right)}_{t}=\left(\frac{Rem}{PIB}\right)_{t}-\overline{\left(\frac{Rem}{PIB}\right)}_{t-1}
\end{equation}

\begin{equation}
\overline{\left(\frac{Rem}{PIB}\right)}_{t-1}=\frac{1}{T}\sum_{j=t-T}^{t-1}\left(\frac{Rem}{PIB}\right)_{j}
\end{equation}


\subsection{Variables endógenas}

\subsubsection{Gobierno}

\subsubsection*{Deuda pública interna}

\begin{equation}
\left(\frac{DPI}{PIB}\right)_{t}=\widetilde{\left(\frac{DPI}{PIB}\right)}_{t}+\overline{\left(\frac{DPI}{PIB}\right)}_{t-1}
\end{equation}

\begin{equation}
\widetilde{\left(\frac{DPI}{PIB}\right)}_{t}=\left(\frac{DPI}{PIB}\right)_{t}-\overline{\left(\frac{DPI}{PIB}\right)}_{t-1}
\end{equation}

\begin{equation}
\overline{\left(\frac{DPI}{PIB}\right)}_{t-1}=\frac{1}{T}\sum_{j=t-T}^{t-1}\left(\frac{DPI}{PIB}\right)_{j}
\end{equation}


\subsubsection{Banco central}

\subsubsection*{Activos externos netos}

\begin{equation}
\left(\frac{AEN}{PIB}\right)_{t}=\widetilde{\left(\frac{AEN}{PIB}\right)}_{t}+\overline{\left(\frac{AEN}{PIB}\right)}_{t-1}
\end{equation}

\begin{equation}
\widetilde{\left(\frac{AEN}{PIB}\right)}_{t}=\left(\frac{AEN}{PIB}\right)_{t}-\overline{\left(\frac{AEN}{PIB}\right)}_{t-1}
\end{equation}

\begin{equation}
\overline{\left(\frac{AEN}{PIB}\right)}_{t-1}=\frac{1}{T}\sum_{j=t-T}^{t-1}\left(\frac{AEN}{PIB}\right)_{j}
\end{equation}


\subsubsection*{Activos internos netos}

\begin{equation}
\left(\frac{AIN}{PIB}\right)_{t}=\widetilde{\left(\frac{AIN}{PIB}\right)}_{t}+\overline{\left(\frac{AIN}{PIB}\right)}_{t-1}
\end{equation}

\begin{equation}
\widetilde{\left(\frac{AIN}{PIB}\right)}_{t}=\left(\frac{AIN}{PIB}\right)_{t}-\overline{\left(\frac{AIN}{PIB}\right)}_{t-1}
\end{equation}

\begin{equation}
\overline{\left(\frac{AEN}{PIB}\right)}_{t-1}=\frac{1}{T}\sum_{j=t-T}^{t-1}\left(\frac{AEN}{PIB}\right)_{j}
\end{equation}


\subsubsection*{Tasa de interés líder de política monetaria}

\begin{equation}
i_{t}=\widetilde{i}_{t}+\overline{i}_{t-1}
\end{equation}

\begin{equation}
\widetilde{i}_{t}=i_{t}-\overline{i}_{t-1}
\end{equation}

\begin{equation}
\overline{i}_{t-1}=\frac{1}{T}\sum_{j=t-T}^{t-1}i_{j}
\end{equation}


\subsubsection{Sector privado}

\subsubsection*{Depreciación nominal}

\begin{equation}
e_{t}=\widetilde{e}_{t}+\overline{e}_{t-1}
\end{equation}

\begin{equation}
\widetilde{e}_{t}=e_{t}-\overline{e}_{t-1}
\end{equation}

\begin{equation}
\overline{e}_{t-1}=\frac{1}{T}\sum_{j=t-T}^{t-1}e_{j}
\end{equation}


\subsubsection*{Prima por riesgo-país}

\begin{equation}
\rho_{t}=\widetilde{\rho}_{t}+\overline{\rho}_{t-1}
\end{equation}

\begin{equation}
\widetilde{\rho}_{t}=\rho_{t}-\overline{\rho}_{t-1}
\end{equation}

\begin{equation}
\overline{\rho}_{t-1}=\frac{1}{T}\sum_{j=t-T}^{t-1}\rho_{j}
\end{equation}


\subsubsection*{Inflación (curva de Phillips)}

\begin{equation}
\pi_{t}=\widetilde{\pi}_{t}+\overline{\pi}_{t-1}
\end{equation}

\begin{equation}
\widetilde{\pi}_{t}=\pi_{t}-\overline{\pi}_{t-1}
\end{equation}

\begin{equation}
\overline{\pi}_{t-1}=\frac{1}{T}\sum_{j=t-T}^{t-1}\pi_{j}
\end{equation}


\subsubsection*{Crecimiento económico (curva IS)}

\begin{equation}
y_{t}=\widetilde{y}_{t}+\overline{y}_{t-1}
\end{equation}

\begin{equation}
\widetilde{y}_{t}=y_{t}-\overline{y}_{t-1}
\end{equation}

\begin{equation}
\overline{y}_{t-1}=\frac{1}{T}\sum_{j=t-T}^{t-1}y_{j}
\end{equation}


\subsubsection*{Inversión privada}

\begin{equation}
\left(\frac{IPriv}{PIB}\right)_{t}=\widetilde{\left(\frac{IPriv}{PIB}\right)}_{t}+\overline{\left(\frac{IPriv}{PIB}\right)}_{t-1}
\end{equation}

\begin{equation}
\widetilde{\left(\frac{IPriv}{PIB}\right)}_{t}=\left(\frac{IPriv}{PIB}\right)_{t}-\overline{\left(\frac{IPriv}{PIB}\right)}_{t-1}
\end{equation}

\begin{equation}
\overline{\left(\frac{IPriv}{PIB}\right)}_{t-1}=\frac{1}{T}\sum_{j=t-T}^{t-1}\left(\frac{IPriv}{PIB}\right)_{j}
\end{equation}


\subsubsection{Otras definiciones}

\subsubsection*{Déficit fiscal}

\begin{equation}
\left(\frac{Def}{PIB}\right)_{t}=\widetilde{\left(\frac{Def}{PIB}\right)}_{t}+\overline{\left(\frac{Def}{PIB}\right)}_{t-1}
\end{equation}

\begin{equation}
\widetilde{\left(\frac{Def}{PIB}\right)}_{t}=\left(\frac{Def}{PIB}\right)_{t}-\overline{\left(\frac{Def}{PIB}\right)}_{t-1}
\end{equation}

\begin{equation}
\overline{\left(\frac{Def}{PIB}\right)}_{t-1}=\frac{1}{T}\sum_{j=t-T}^{t-1}\left(\frac{Def}{PIB}\right)_{j}
\end{equation}


\subsubsection*{Deuda pública total}

\begin{equation}
\left(\frac{DPT}{PIB}\right)_{t}=\widetilde{\left(\frac{DPT}{PIB}\right)}_{t}+\overline{\left(\frac{DPT}{PIB}\right)}_{t-1}
\end{equation}

\begin{equation}
\widetilde{\left(\frac{DPT}{PIB}\right)}_{t}=\left(\frac{DPT}{PIB}\right)_{t}-\overline{\left(\frac{DPT}{PIB}\right)}_{t-1}
\end{equation}

\begin{equation}
\overline{\left(\frac{DPT}{PIB}\right)}_{t-1}=\frac{1}{T}\sum_{j=t-T}^{t-1}\left(\frac{DPT}{PIB}\right)_{j}
\end{equation}


\subsubsection*{Base monetaria}

\begin{equation}
\left(\frac{MO}{PIB}\right)_{t}=\widetilde{\left(\frac{MO}{PIB}\right)}_{t}+\overline{\left(\frac{MO}{PIB}\right)}_{t-1}
\end{equation}

\begin{equation}
\widetilde{\left(\frac{MO}{PIB}\right)}_{t}=\left(\frac{MO}{PIB}\right)_{t}-\overline{\left(\frac{MO}{PIB}\right)}_{t-1}
\end{equation}

\begin{equation}
\overline{\left(\frac{MO}{PIB}\right)}_{t-1}=\frac{1}{T}\sum_{j=t-T}^{t-1}\left(\frac{MO}{PIB}\right)_{j}
\end{equation}


\section{Algunas convenciones}

En esta sección se presentan algunas convenciones adoptadas respecto
al tratamiento de variables de flujo, stock, tasas de interés y tasas
de variación.

\subsection*{Variables flujo}

Este tipo de variables se tratan como la suma móvil de 4 trimestres.

\begin{equation}
x_{t}=\sum_{j=t-4+1}^{t}x_{Q,j}
\end{equation}


\subsection*{Variables stock}

Las variables stock son un promedio móvil de 4 trmestres.

\begin{equation}
\bar{x}_{t}=\frac{1}{4}\sum_{j=t-4+1}^{t}x_{Q,j}
\end{equation}


\subsection*{Tasas de interés}

Las tasas de interés doméstica y externa están definidas como tasas
anuales.

\subsection*{Tasas de crecimiento económico}

Esta variable se define como la variación interanual de la suma móvil
de 4 trimestres del PIB.

\section{Simplificación notacional}

Para simplificar la notación del modelo, se introducen las definiciones
siguientes:

\begin{equation}
y_{1,t}=\widetilde{y}_{t}^{*}
\end{equation}

\begin{equation}
y_{2,t}=\widetilde{i}_{t}^{*}
\end{equation}

\begin{equation}
y_{3,t}=\widetilde{\pi}_{t}^{*}
\end{equation}

\begin{equation}
y_{4,t}=\widetilde{\left(\frac{IPub}{PIB}\right)}_{t}
\end{equation}

\begin{equation}
y_{5,t}=\widetilde{\left(\frac{IngT}{PIB}\right)}_{t}
\end{equation}

\begin{equation}
y_{6,t}=\widetilde{\left(\frac{GFun}{PIB}\right)}_{t}
\end{equation}

\begin{equation}
y_{7,t}=\widetilde{\left(\frac{DPE}{PIB}\right)}_{t}
\end{equation}

\begin{equation}
y_{8,t}=\widetilde{\left(\frac{Rem}{PIB}\right)}_{t}
\end{equation}

\begin{equation}
y_{9,t}=\widetilde{\left(\frac{DPI}{PIB}\right)}_{t}
\end{equation}

\begin{equation}
y_{10,t}=\widetilde{\left(\frac{AEN}{PIB}\right)}_{t}
\end{equation}

\begin{equation}
y_{11,t}=\widetilde{\left(\frac{AIN}{PIB}\right)}_{t}
\end{equation}

\begin{equation}
y_{12,t}=\widetilde{i}_{t}
\end{equation}

\begin{equation}
y_{13,t}=\widetilde{e}_{t}
\end{equation}

\begin{equation}
y_{14,t}=\widetilde{\rho}_{t}
\end{equation}

\begin{equation}
y_{15,t}=\widetilde{\pi}_{t}
\end{equation}

\begin{equation}
y_{16,t}=\widetilde{y}_{t}
\end{equation}

\begin{equation}
y_{17,t}=\widetilde{\left(\frac{IPriv}{PIB}\right)}_{t}
\end{equation}


\section{Representación compacta y simplificada del modelo}

Sustituyendo las ecuaciones (121) a (137) en el modelo compacto (ecuaciones
(38) a (55)), se obtiene la versión compacta y simplificada del modelo,
como sigue:

\begin{equation}
y_{1,t}=\alpha_{1,1}y_{1,t-1}+u_{1,t}
\end{equation}

\begin{equation}
y_{2,t}=\alpha_{2,1}y_{2,t-1}+u_{2,t}
\end{equation}

\begin{equation}
y_{3,t}=\alpha_{3,1}y_{3,t-1}+u_{3,t}
\end{equation}

\begin{equation}
y_{4,t}=\alpha_{4,1}y_{4,t-1}+u_{4,t}
\end{equation}

\begin{equation}
y_{5,t}=\alpha_{5,1}y_{5,t-1}+u_{5,t}
\end{equation}

\begin{equation}
y_{6,t}=\alpha_{6,1}y_{6,t-1}+u_{6,t}
\end{equation}

\begin{equation}
y_{7,t}=\alpha_{7,1}y_{7,t-1}+u_{7,t}
\end{equation}

\begin{equation}
y_{8,t}=\alpha_{8,1}y_{8,t-1}+u_{8,t}
\end{equation}

\begin{equation}
y_{9,t}=\alpha_{9,1}y_{9,t-1}+\alpha_{9,2}y_{7,t-1}+\alpha_{9,3}\left(y_{6,t-1}+y_{4,t-1}-y_{5,t-1}\right)+u_{9,t}
\end{equation}

\begin{equation}
y_{10,t}=\alpha_{10,1}y_{10,t-1}+\alpha_{10,2}y_{7,t-1}+\alpha_{10,3}y_{3,t-1}+\alpha_{10,4}y_{13,t-1}+\alpha_{10,5}y_{15,t-1}+u_{10,t}
\end{equation}

\begin{equation}
y_{11,t}=\alpha_{11,1}y_{11,t-1}+\alpha_{11,2}y_{9,t-1}+\alpha_{11,3}y_{12,t-1}+\alpha_{11,4}y_{15,t-1}+u_{11,t}
\end{equation}

\begin{equation}
y_{12,t}=\alpha_{12,1}y_{12,t-1}+\alpha_{12,2}y_{2,t-1}+\alpha_{12,3}y_{15,t-1}+\alpha_{12,4}y_{16,t-1}+u_{12,t}
\end{equation}

\begin{equation}
y_{13,t}=\alpha_{13,1}y_{13,t-1}+\alpha_{13,2}\left(y_{7,t-1}-y_{10,t-1}\right)+\alpha_{13,3}\left(y_{12,t-1}-y_{2,t-1}-y_{14,t-1}\right)+u_{13,t}
\end{equation}

\begin{equation}
y_{14,t}=\alpha_{14,1}y_{14,t-1}+\alpha_{14,2}\left(y_{9,t-1}-y_{11,t-1}+y_{7,t-1}-y_{10,t-1}\right)+u_{14,t}
\end{equation}

\begin{equation}
y_{15,t}=\alpha_{15,1}y_{15,t-1}+\alpha_{15,2}y_{16,t-1}+\alpha_{15,3}y_{3,t-1}+\alpha_{15,4}y_{13,t-1}+\alpha_{15,5}\left(y_{10,t-1}+y_{11,t-1}\right)+u_{15,t}
\end{equation}

\begin{equation}
y_{16,t}=\alpha_{16,1}y_{16,t-1}+\alpha_{16,2}y_{1,t-1}+\alpha_{16,3}y_{8,t-1}+\alpha_{16,4}\left(y_{4,t-1}+y_{17,t-1}\right)+\alpha_{16,5}\left(y_{10,t-1}+y_{11,t-1}\right)+u_{16,t}
\end{equation}

\begin{equation}
y_{17,t}=\alpha_{17,1}y_{17,t-1}+\alpha_{17,2}y_{12,t-1}+\alpha_{17,3}\left(y_{9,t-1}-y_{11,t-1}\right)+\alpha_{17,4}y_{14,t-1}+u_{17,t}
\end{equation}

\[
u_{9,t}=\left(1-\alpha_{9,1}\right)y_{9,t-1}+\left(1-\alpha_{7,1}-\alpha_{9,2}\right)y_{7,t-1}+\left(\alpha_{6,1}-\alpha_{9,3}\right)y_{6,t-1}
\]

\begin{equation}
+\left(\alpha_{4,1}-\alpha_{9,3}\right)y_{4,t-1}-\left(\alpha_{5,1}+\alpha_{9,3}\right)y_{5,t-1}+u_{4,t}+u_{6,t}-u_{5,t}-u_{7,t}
\end{equation}


\section{Representación matricial del modelo}

Partiendo de las ecuaciones presentadas en la sección 9, a continuación
se representa el modelo en forma matricial.

\subsection*{Matriz de variables del modelo en forma compacta y simplificada}

Donde $y_{t}$ es una matriz de 17$\times$1, conteniendo las variables
definidas en la sección anterior.

\begin{equation}
Y_{t}=\begin{bmatrix}y_{1,t}\\
y_{2,t}\\
\vdots\\
y_{4,t}
\end{bmatrix}
\end{equation}


\subsection*{Matriz de coeficientes}

La matriz $A_{1}$ es de dimensiones 17$\times$17 y contiene los
coeficientes del modelo. Es importante hacer notar que el modelo tiene
expresiones del tipo $\alpha_{9,3}\left(y_{6,t-1}+y_{4,t-1}-y_{5,t-1}\right)$
(ecuación 145) e indica que las variables $y_{6,t}$, $y_{4,t}$,
y $y_{5,t}$ comparten el mismo coeficiente y en el caso de $y_{5,t}$
este se encuentra con el signo negativo. Este tipo de expresiones
se puede representar de forma apropiada en la matriz de coeficientes
mostrada a continuación: 
\begin{flushleft}
{\tiny
\begin{equation}
\hspace*{-2cm}A_{1}=\setlength{\arraycolsep}{3pt}\begin{bmatrix}\alpha_{1,1} & 0 & 0 & 0 & 0 & 0 & 0 & 0 & 0 & 0 & 0 & 0 & 0 & 0 & 0 & 0 & 0\\
0 & \alpha_{2,1} & 0 & 0 & 0 & 0 & 0 & 0 & 0 & 0 & 0 & 0 & 0 & 0 & 0 & 0 & 0\\
0 & 0 & \alpha_{3,1} & 0 & 0 & 0 & 0 & 0 & 0 & 0 & 0 & 0 & 0 & 0 & 0 & 0 & 0\\
0 & 0 & 0 & \alpha_{4,1} & 0 & 0 & 0 & 0 & 0 & 0 & 0 & 0 & 0 & 0 & 0 & 0 & 0\\
0 & 0 & 0 & 0 & \alpha_{5,1} & 0 & 0 & 0 & 0 & 0 & 0 & 0 & 0 & 0 & 0 & 0 & 0\\
0 & 0 & 0 & 0 & 0 & \alpha_{6,1} & 0 & 0 & 0 & 0 & 0 & 0 & 0 & 0 & 0 & 0 & 0\\
0 & 0 & 0 & 0 & 0 & 0 & \alpha_{7,1} & 0 & 0 & 0 & 0 & 0 & 0 & 0 & 0 & 0 & 0\\
0 & 0 & 0 & 0 & 0 & 0 & 0 & \alpha_{8,1} & 0 & 0 & 0 & 0 & 0 & 0 & 0 & 0 & 0\\
0 & 0 & 0 & \alpha_{9,3} & -\alpha_{9,3} & \alpha_{9,3} & \alpha_{9,2} & 0 & \alpha_{9,1} & 0 & 0 & 0 & 0 & 0 & 0 & 0 & 0\\
0 & 0 & \alpha_{10,3} & 0 & 0 & 0 & \alpha_{10,2} & 0 & 0 & \alpha_{10,1} & 0 & 0 & 0 & 0 & \alpha_{10,5} & 0 & 0\\
0 & 0 & 0 & 0 & 0 & 0 & 0 & 0 & \alpha_{11,2} & 0 & \alpha_{11,1} & \alpha_{11.3} & 0 & 0 & \alpha_{11,4} & 0 & 0\\
0 & \alpha_{12,2} & 0 & 0 & 0 & 0 & 0 & 0 & 0 & 0 & 0 & \alpha_{12,1} & 0 & 0 & \alpha_{12,3} & \alpha_{12,4} & 0\\
0 & -\alpha_{13,3} & 0 & 0 & 0 & 0 & \alpha_{13,2} & 0 & 0 & \alpha_{13,2} & 0 & \alpha_{13,3} & \alpha_{13,1} & -\alpha_{13,3} & 0 & 0 & 0\\
0 & 0 & 0 & 0 & 0 & 0 & \alpha_{14,2} & 0 & \alpha_{14,2} & -\alpha_{14,2} & -\alpha_{14,2} & 0 & 0 & \alpha_{14,1} & 0 & 0 & 0\\
0 & 0 & \alpha_{15,3} & 0 & 0 & 0 & 0 & 0 & 0 & \alpha_{15,5} & \alpha_{15,5} & 0 & 0 & 0 & \alpha_{15,1} & \alpha_{15,2} & 0\\
\alpha_{16,1} & 0 & 0 & \alpha_{16,4} & 0 & 0 & 0 & \alpha_{16,3} & 0 & \alpha_{16,5} & \alpha_{16,5} & 0 & 0 & 0 & 0 & \alpha_{16,1} & \alpha_{16,4}\\
0 & 0 & 0 & 0 & 0 & 0 & 0 & 0 & \alpha_{17,3} & 0 & -\alpha_{17,3} & \alpha_{17,2} & 0 & \alpha_{17,4} & 0 & 0 & \alpha_{17,1}
\end{bmatrix}
\end{equation}
}{\tiny\par}
\par\end{flushleft}

\subsection*{Matriz de residuos}

La matriz $u_{t}$es de dimensiones 17$\times$1 y contiene los residuos
del modelo.

\begin{equation}
U_{t}=\begin{bmatrix}u_{1,t}\\
u_{2,t}\\
\vdots\\
u_{17,t}
\end{bmatrix}
\end{equation}

Finalmente el modelo se presenta mediante de forma matricial de la
siguiente manera:

\begin{equation}
Y_{t}=AY_{t-1}+U_{t}
\end{equation}

\end{document}
